% =========================================================
% The Six Exercise Types in Analysis
% =========================================================
\paragraph{The Six Exercise Types.}

Every exercise in sequences and metric spaces belongs to one of six
structural types. Identifying the type is the first step before
writing anything. The type determines the shape of the solution ---
not the details, but the overall structure.

\begin{tcolorbox}[colback=gray!6, colframe=gray!40, arc=2pt,
  left=6pt, right=6pt, top=4pt, bottom=4pt,
  title={\small\textbf{Quick Reference: The Six Types}},
  fonttitle=\small\bfseries]
\small
\renewcommand{\arraystretch}{1.3}
\begin{tabular}{p{0.13\textwidth} p{0.28\textwidth} p{0.25\textwidth} p{0.20\textwidth}}
\toprule
\textbf{Type} & \textbf{Signal words} & \textbf{Shape} & \textbf{First move} \\
\midrule
Proof &
  \textit{Prove, show, establish} &
  Hypotheses to conclusion &
  Unpack all definitions \\
Verification &
  \textit{Verify that \ldots is a metric, show \ldots satisfies} &
  Check each condition &
  List every condition \\
Prove/Counterexample &
  \textit{True? Prove or find a counterexample} &
  Decide, then justify &
  Decide first \\
Construction &
  \textit{Find, construct, exhibit, give an example of} &
  Name the object; verify it &
  Name the object \\
Computation &
  \textit{Compute, calculate, find the diameter of} &
  Evaluate from a definition &
  Write the definition \\
Characterization &
  \textit{If and only if, characterize all \ldots that} &
  Two separate directions &
  Split two directions \\
\bottomrule
\end{tabular}
\end{tcolorbox}

\begin{remark}[Type 1 (Proof): you are told the statement is true]
In a proof exercise, the statement is true and your job is to trace
the logical path from the hypotheses to the conclusion. There is no
choice to make about truth value.

The opening move is always the same: unpack every definition in the
hypotheses and in the conclusion. The definitions turn abstract concepts
into concrete quantifier-and-inequality statements that can be
manipulated. Without this step, you have nothing to work with.

\textbf{Example.} \textit{Prove that every convergent sequence is Cauchy.}
Unpack ``convergent'': $\forall\,\varepsilon>0\;\exists\, N$ s.t.\
$n \geq N \Rightarrow d(x_n, p) < \varepsilon$.
Unpack ``Cauchy'': $\forall\,\varepsilon>0\;\exists\, M$ s.t.\
$m,n \geq M \Rightarrow d(x_m, x_n) < \varepsilon$.
Now the proof goal is visible: use the convergence condition (which
controls $d(x_n, p)$) to establish the Cauchy condition (which asks
for control over $d(x_m, x_n)$). The triangle inequality bridges them.
\end{remark}

\begin{remark}[Type 2 (Verification): check every condition]
In a verification exercise, a specific named object is given and you
must confirm it satisfies every condition of a definition. The definition
has a finite list of conditions. Your job is to work through every one,
explicitly and separately.

The error students make: checking four out of five conditions and
writing QED. A verification is complete only when every condition
is checked. When you begin a verification, write the complete list of
conditions first. Then work through them in order.

\textbf{Example.} \textit{Verify that the discrete metric is a metric.}
A metric has four conditions: non-negativity, identity of indiscernibles,
symmetry, triangle inequality. List all four. Check each one.
The triangle inequality requires a case split (on whether $x = z$ or
$x \neq z$), which is itself a standard move in analysis.
\end{remark}

\begin{remark}[Type 3 (Prove or counterexample): decide before you argue]
This is the most dangerous type for students who rush. The exercise
gives you a statement and asks you to determine its truth value, then
justify. If you skip the determination stage and try to prove a false
statement, you will waste effort and grow confused.

Stage~1 is always: decide. Test the statement on simple cases ---
$\mathbb{R}$ with the standard metric first, then discrete spaces,
then other special cases. One counterexample falsifies; failure to
find a counterexample after testing several cases suggests the
statement is true.

Stage~2: justify. If true, write a proof. If false, exhibit a specific
counterexample and verify it fails.

\textbf{Example.} \textit{Is the open ball always equal to the closed ball?}
Test in $\mathbb{R}$: $B(0,1) = (-1,1)$ and $C(0,1) = [-1,1]$.
The point $-1$ is in $C(0,1)$ but not $B(0,1)$. False. Exhibit this
specific counterexample in the solution, with the witness ($-1$) and
the calculation ($d(-1,0) = 1$, so $-1 \in C(0,1)$ and $-1 \notin B(0,1)$).
\end{remark}

\begin{remark}[Type 4 (Construction): name first, verify second]
A construction exercise asks you to produce an object with given
properties. The solution has two parts: stating the object, and
verifying it has every required property. Both parts are required.

State the object completely and unambiguously before any verification.
If you cannot state it clearly, you do not yet have a solution.
Then verify every required property from scratch --- do not assume
properties that were not established.

\textbf{Example.} \textit{Find a metric on $\mathbb{N}$ such that
$\mathbb{N}$ is bounded.} The idea: embed $\mathbb{N}$ into
$(0,1] \subset \mathbb{R}$ via $n \mapsto 1/n$ and pull back the metric.
Define $d(m,n) = |1/m - 1/n|$. Then verify: (i) it is a metric (four
conditions), and (ii) it is bounded ($d(m,n) \leq 1$ for all $m,n$).
\end{remark}

\begin{remark}[Type 5 (Computation): the definition is the answer]
In a computation exercise, a specific numerical or set-theoretic value
is asked for. The answer always comes from writing out the relevant
definition as a supremum, infimum, or set, and then evaluating it.

Do not skip straight to the answer. Write the definition first.
This is not just for clarity --- it prevents errors that arise from
misremembering what is being computed.

\textbf{Example.} \textit{Compute $\operatorname{diam}([a,b])$.}
Write: $\operatorname{diam}([a,b]) = \sup\{|x-y| : x, y \in [a,b]\}$.
Upper bound: $|x-y| \leq b-a$ for all $x,y \in [a,b]$.
Attained: $|b - a| = b-a$. Therefore $\operatorname{diam}([a,b]) = b-a$.
\end{remark}

\begin{remark}[Type 6 (Characterization): two proofs, not one]
A characterization exercise asks for an equivalence or an if-and-only-if.
The solution is always two separate proofs, labeled $(\Rightarrow)$ and
$(\Leftarrow)$. Label them before writing a single line.

The two directions often require completely different techniques and
different amounts of work. The $(\Rightarrow)$ direction may be three
lines while the $(\Leftarrow)$ direction requires a sequence argument.
Writing them separately prevents the error of running arguments in
both directions simultaneously and losing track of which assumptions
are in play.

\textbf{Example.} \textit{Show that a sequence converges if and only if
it is eventually constant (in the discrete metric).}
$(\Leftarrow)$: Suppose eventually constant; then convergence holds trivially.
$(\Rightarrow)$: Suppose it converges; apply the definition with
$\varepsilon = 1/2$ and use the fact that distances in the discrete metric
are either $0$ or $1$.
\end{remark}
